% ============================================================================================
% BAB III ANALISIS MASALAH
% Pembagian subbab tidak rigid dan dapat bervariasi. Bab ini minimal berisi analisis kebutuhan
% fungsional dan nonfungsional, analisis berbagai alternatif solusi yang dapat ditawarkan, dan
% metode pemilihan solusi yang diusulkan.
% ============================================================================================

\chapter{ANALISIS MASALAH}
\label{chap:analisis-masalah}
\section{Analisis Kondisi Saat Ini}
\label{sec:kondisi_saat_ini}

% Analisis terhadap lanskap teknologi pembelajaran adaptif saat ini mengungkapkan sebuah paradoks fundamental: sementara teknologi mitigasi \textit{over-personalization} telah matang di domain sistem rekomendasi komersial, adopsinya di domain pendidikan terhambat oleh arsitektur sistem asesmen yang kaku.

% \subsection{Keterbatasan Arsitektur Asesmen dalam Mendeteksi Isolasi}
Masalah utama saat ini bukan sekadar ketiadaan fitur sosial, melainkan ketidakmampuan algoritma inti (\textit{core algorithms}) untuk memproses sinyal sosial. Teknik dominan seperti \textit{Bayesian Knowledge Tracing} (BKT) atau \textit{Elo-rating} bekerja dengan asumsi dasar bahwa pembelajaran adalah proses soliter \autocite{Minn2022}. Variabel yang diolah hanyalah interaksi biner antara siswa dan konten (benar/salah) \autocite{Vesin2022}.

Implikasinya adalah sistem saat ini menderita \textit{algorithmic blindness}
terhadap isolasi pembelajaran. Meskipun seorang siswa terjebak dalam
\textit{filter bubble} atau \textit{echo chamber} akademik, hanya mengerjakan
soal yang ia sukai dan tidak pernah berinteraksi dengan perspektif lain, 
sistem
asesmen konvensional tetap akan menilainya berkinerja tinggi selama ia menjawab
benar. Kondisi ini berbahaya karena menciptakan "kompetensi semu", suatu fenomena ketika siswa
mahir secara individu namun rapuh dalam kolaborasi, yang bertentangan dengan visi
ekosistem kolaboratif \autocite{BPPT2020}.

% \subsection{Analisis Kesenjangan Adopsi Solusi Mitigasi \textit{Over-Personalization}}
Studi literatur menunjukkan bahwa solusi untuk masalah isolasi ini sudah tersedia
di bidang lain, seperti \textit{User-Controllable Recommender Systems} (UCRS)
\autocite{Wang2022} dan \textit{Group Recommender Systems} (GRS)
\autocite{Dokoupil2022}. Namun, analisis kondisi saat ini menunjukkan bahwa
solusi-solusi tersebut gagal diterapkan secara efektif dalam pendidikan karena
alasan arsitektural berikut:

\begin{enumerate}
    \item \textbf{Kegagalan Pendekatan Kendali Pengguna (UCRS):} Solusi \textcite{Wang2022} mengandalkan pengguna untuk sadar bahwa mereka terjebak dalam \textit{bubble}. Dalam konteks pendidikan, siswa (sebagai \textit{novice}) sering kali tidak memiliki metakognisi yang cukup untuk menyadari bahwa mereka butuh materi yang lebih beragam atau menantang secara sosial. Oleh karena itu, intervensi manual ala UCRS tidak cukup; diperlukan intervensi otomatis dari sistem.
    
    \item \textbf{Kegagalan Pendekatan Visual (OSSM):} Upaya yang ada saat ini seperti \textit{Open Social Student Modeling} (OSSM) hanya menempatkan data sosial sebagai "lapisan visual" di atas sistem, tanpa mengubah logika rekomendasi di bawahnya \autocite{Brusilovsky2016}. Data perbandingan sosial hanya menjadi "hiasan" motivasi, tidak menjadi variabel input yang mengubah tingkat kesulitan atau jenis materi yang direkomendasikan algoritma.
\end{enumerate}

% \subsection{Implikasi Terhadap Arah Pengembangan Sistem (\textit{What's Next})}
Dari analisis di atas, dapat disimpulkan bahwa menambahkan fitur kolaborasi saja (seperti forum atau grup) tidak akan memecahkan masalah \textit{over-personalization} jika mesin asesmennya tetap terpisah.

Kebutuhan mendesak saat ini adalah melakukan rekayasa ulang (\textit{re-engineering}) pada alur data. Kita tidak lagi membutuhkan dua sistem yang berjalan paralel (satu untuk nilai individu, satu untuk grup), melainkan sebuah sistem hibrida. Data kinerja sosial, seperti kualitas umpan balik dalam \textit{peer assessment} \autocite{Lu2012}, harus diubah menjadi variabel matematis yang dapat "dibaca" oleh algoritma asesmen adaptif (seperti parameter $K$ dalam \textit{Elo-rating}).

Hanya dengan cara inilah, sistem dapat secara proaktif mendeteksi isolasi dan melakukan "koreksi paksa" melalui rekomendasi konten, menjembatani kesenjangan menuju analisis kebutuhan sistem yang akan diuraikan pada subbab berikutnya.

\section{Analisis Kebutuhan}
\label{sec:analisis_kebutuhan}

Berdasarkan analisis kondisi saat ini pada Subbab \ref{sec:kondisi_saat_ini}, sistem yang dibutuhkan harus mampu menjembatani keterpisahan antara personalisasi individu dengan lingkungan sosialnya untuk memitigasi risiko \textit{over-personalization}. Kebutuhan ini dirumuskan tanpa membatasi teknik implementasi spesifik, guna memungkinkan eksplorasi berbagai alternatif solusi (seperti \textit{peer assessment}, \textit{mentoring}, atau rekomendasi jejaring) pada tahap pemilihan solusi.

\subsection{Kebutuhan Fungsional}
\label{ssec:kebutuhan_fungsional}
Daftar kebutuhan fungsional sistem dari sisi pengguna (siswa) dan logika sistem disajikan pada Tabel \ref{tbl:kebutuhan-fungsional}.

\begin{longtable}{@{\extracolsep{\fill}} l p{0.8\textwidth}}
\caption{Kebutuhan Fungsional (FR) Sistem}\label{tbl:kebutuhan-fungsional} \\
\toprule
\textbf{Kode} & \textbf{Deskripsi Kebutuhan Fungsional} \\
\midrule
\endfirsthead

\caption*{Tabel \ref{tbl:kebutuhan-fungsional} Kebutuhan Fungsional (FR) Sistem (lanjutan)} \\
\toprule
\textbf{Kode} & \textbf{Deskripsi Kebutuhan Fungsional} \\
\midrule
\endhead

\midrule
\multicolumn{2}{r}{\textit{Bersambung ke halaman berikutnya}} \\
\bottomrule
\endfoot

\bottomrule
\endlastfoot

% === Data Tabel Kebutuhan Fungsional ===
\multicolumn{2}{l}{\textbf{Interaksi Pengguna (Siswa)}} \\
\addlinespace[0.3em]
FR-01 & Sistem harus menyediakan mekanisme otentikasi bagi siswa untuk mengakses lingkungan pembelajaran personal mereka. \\
FR-02 & Sistem harus menyajikan materi pembelajaran yang telah disesuaikan tingkat kesulitannya dengan kemahiran siswa saat ini. \\
FR-03 & Sistem harus memfasilitasi \textbf{interaksi pembelajaran antar-pengguna} yang memungkinkan terjadinya pertukaran informasi atau kolaborasi akademik di dalam platform. \\
FR-04 & Sistem harus menyediakan informasi mengenai konteks sosial pembelajaran (misalnya, aktivitas komunitas atau capaian kolektif) untuk membangun kesadaran sosial (\textit{social awareness}). \\
\addlinespace[0.8em]

\multicolumn{2}{l}{\textbf{Logika Adaptasi \& Pemrosesan Data (Sistem)}} \\
\addlinespace[0.3em]
FR-05 & Sistem harus memiliki algoritma untuk mengestimasi tingkat kemahiran pengetahuan individu siswa secara dinamis berdasarkan riwayat aktivitasnya. \\
FR-06 & Sistem harus mampu \textbf{mengkuantifikasi data interaksi sosial} (dari FR-03) menjadi parameter numerik yang dapat diproses lebih lanjut (misalnya, skor keterlibatan atau skor kinerja kolaboratif). \\
FR-07 & Sistem harus memiliki mekanisme untuk menggabungkan parameter kemahiran individu (FR-05) dan parameter sosial (FR-06) dalam satu model komputasi terpadu. \\
FR-08 & Sistem harus menggunakan hasil penggabungan model (FR-07) untuk menentukan materi atau aktivitas berikutnya yang paling optimal bagi siswa. \\
FR-09 & Sistem harus mampu merekomendasikan \textbf{intervensi sosial spesifik} (misalnya, merekomendasikan orang untuk berinteraksi atau topik untuk didiskusikan) guna mencegah isolasi pembelajaran. \\
\end{longtable}

\subsection{Kebutuhan Nonfungsional}
\label{ssec:kebutuhan_nonfungsional}
Kebutuhan nonfungsional disesuaikan untuk mendukung interaksi sosial yang responsif tanpa menentukan metode spesifik.

\begin{longtable}{@{\extracolsep{\fill}} l p{0.8\textwidth}}
\caption{Kebutuhan Nonfungsional (NFR) Sistem}\label{tbl:kebutuhan-nonfungsional} \\
\toprule
\textbf{Kode} & \textbf{Deskripsi Kebutuhan Nonfungsional} \\
\midrule
\endfirsthead

\caption*{Tabel \ref{tbl:kebutuhan-nonfungsional} Kebutuhan Nonfungsional (NFR) Sistem (lanjutan)} \\
\toprule
\textbf{Kode} & \textbf{Deskripsi Kebutuhan Nonfungsional} \\
\midrule
\endhead

\midrule
\multicolumn{2}{r}{\textit{Bersambung ke halaman berikutnya}} \\
\bottomrule
\endfoot

\bottomrule
\endlastfoot

% === Data Tabel Kebutuhan Nonfungsional ===
NFR-01 & \textbf{(Usability)} Antarmuka interaksi sosial harus terintegrasi secara intuitif dengan materi pembelajaran, meminimalkan beban kognitif siswa saat berpindah antara belajar mandiri dan berkolaborasi. \\
NFR-02 & \textbf{(Interpretability)} Sistem harus mampu memberikan indikasi alasan di balik rekomendasi (baik materi maupun interaksi sosial) agar dapat dipahami oleh siswa. \\
NFR-03 & \textbf{(Performance)} Pembaruan model pengguna (baik skor individu maupun sosial) harus terjadi dalam waktu nyata (\textit{near real-time}) setelah interaksi selesai dilakukan. \\
NFR-04 & \textbf{(Scalability)} Sistem harus mampu menangani konkurensi interaksi sosial dari satu kelas (±50 pengguna) secara simultan. \\
NFR-05 & \textbf{(Privacy)} Mekanisme interaksi sosial harus dirancang sedemikian rupa sehingga data privasi spesifik (seperti nilai ujian mentah) tidak terekspos ke pengguna lain tanpa izin. \\
NFR-06 & \textbf{(Access)} Sistem dapat diakses melalui peramban web standar. \\
\end{longtable}

\subsection{Pemetaan Kebutuhan}
\label{ssec:pemetaan_kebutuhan}
Tabel \ref{tbl:pemetaan-kebutuhan} memetakan tujuan penanganan masalah ke fungsi spesifik yang telah digeneralisasi.

\begin{table}[H]
    \centering
    \captionsetup{justification=centering}
    \caption{Pemetaan Masalah dan Kebutuhan Fungsional}
    \label{tbl:pemetaan-kebutuhan}
    \begin{tabular}{@{\extracolsep{\fill}} l p{0.55\textwidth}}
    \toprule
    \textbf{Fokus Penanganan Masalah} & \textbf{Implementasi Fungsional (FR)} \\
    \midrule
    Model Kompetensi Individu & FR-02, FR-05 \\
    Fasilitasi Interaksi Sosial & FR-01, FR-03, FR-04 \\
    Kuantifikasi Aspek Sosial & FR-06 \\
    Integrasi Model & FR-07 \\
    Adaptasi & FR-08, FR-09 \\
    \bottomrule
    \end{tabular}
\end{table}

\section{Analisis Pemilihan Solusi}
\subsection{Alternatif Solusi}
\label{ssec:alternatif_solusi}
Untuk menjawab kebutuhan sistem dalam memitigasi \textit{over-personalization}, diperlukan pemilihan komponen teknologi yang tidak hanya akurat, tetapi juga fleksibel untuk diintegrasikan. Pemilihan solusi dilakukan pada tiga dimensi utama: (1) teknik asesmen kemahiran individu, (2) teknik interaksi dan asesmen sosial, dan (3) teknik spesifik pencegahan \textit{echo chamber}.

\subsubsection{Alternatif Teknik Asesmen Kemahiran Individu (R1)}
\label{sssec:alternatif_asesmen_individu}
Komponen ini berfungsi memodelkan kompetensi siswa. Dalam konteks pencegahan \textit{over-personalization}, kriteria utamanya adalah apakah model tersebut cukup fleksibel untuk "diintervensi" oleh parameter sosial tanpa merusak integritas matematisnya.

\begin{itemize}
    \item \textbf{\textit{Item Response Theory} (IRT) dan \textit{Bayesian IRT} (IRT*):}
    Pendekatan psikometri standar yang memodelkan probabilitas jawaban benar berdasarkan parameter siswa ($\theta$) dan butir soal ($b$) \autocite{Handayani2025}.
    \textbf{Analisis terhadap \textit{Over-personalization}:} Model IRT cenderung statis dan mengasumsikan parameter invarian. Sangat sulit untuk menyuntikkan "bobot sosial" ke dalam rumus IRT secara dinamis untuk mengubah rekomendasi soal secara \textit{real-time} saat terdeteksi adanya isolasi pembelajaran \autocite{Vesin2022}.

    \item \textbf{\textit{Factor Analysis} (FA):}
    Teknik statistik yang digunakan untuk mengidentifikasi struktur faktor laten (seperti pengetahuan awal atau kemampuan kognitif) yang mendasari kinerja siswa \autocite{Ihichr2024}.
    \textbf{Analisis terhadap \textit{Over-personalization}:} FA lebih cocok digunakan untuk validasi struktur tes atau analisis \textit{post-hoc} (setelah data terkumpul) daripada pelacakan dinamis. Teknik ini tidak dirancang untuk pembaruan parameter \textit{real-time} yang diperlukan untuk merespons perubahan perilaku sosial siswa seketika.

    \item \textbf{\textit{Bayesian Knowledge Tracing} (BKT):}
    Model berbasis \textit{Hidden Markov Models} (HMM) yang melacak transisi biner status pengetahuan dari "belum belajar" ke "sudah belajar" \autocite{Minn2022}.
    \textbf{Analisis terhadap \textit{Over-personalization}:} BKT memiliki struktur probabilitas yang kaku (mengandalkan parameter \textit{guess} dan \textit{slip}). Memasukkan variabel sosial kontinu ke dalam probabilitas transisi BKT sangat kompleks dan berisiko melanggar asumsi dasar model bahwa keterampilan tidak akan "dilupakan" \autocite{Ihichr2024}, padahal isolasi sosial mungkin memerlukan penurunan skor untuk memicu intervensi.

    \item \textbf{\textit{Deep Knowledge Tracing} (DKT):}
    Menggunakan jaringan saraf tiruan (\textit{Recurrent Neural Networks}) untuk melacak pengetahuan \autocite{Minn2022}.
    \textbf{Analisis terhadap \textit{Over-personalization}:} Meskipun akurasinya tinggi, sifatnya yang \textit{black-box} membuat sistem sulit mendeteksi \textit{kenapa} seorang siswa terjebak dalam pola soal tertentu. Intervensi algoritmik untuk memecahkan \textit{bubble} sulit dilakukan secara terukur karena kompleksitas parameternya.

    \item \textbf{Sistem \textit{Elo-rating}:}
    Memodelkan interaksi siswa-konten sebagai kompetisi dengan rumus pembaruan sederhana: $R_{baru} = R_{lama} + K(S - E)$ \autocite{Vesin2022}.
    \textbf{Analisis terhadap \textit{Over-personalization}:} Ini adalah kandidat terkuat karena keberadaan faktor $K$ (koefisien bobot). Faktor ini dapat dimodifikasi secara eksplisit menjadi fungsi dari aktivitas sosial (misalnya, siswa yang kurang bersosialisasi mendapat $K$ yang memaksa perubahan peringkat lebih cepat), memungkinkan integrasi langsung untuk melawan stagnasi algoritma.
\end{itemize}

\subsubsection{Alternatif Teknik Interaksi dan Asesmen Sosial (R2 \& R3)}
\label{sssec:alternatif_asesmen_sosial}
Komponen ini bertujuan memfasilitasi interaksi yang bermakna. Untuk memecahkan \textit{filter bubble}, aktivitas sosial harus memaksa siswa terpapar pada perspektif atau materi di luar preferensi historis mereka.

\begin{itemize}
    \item \textbf{Pembelajaran Berbasis Proyek (\textit{Project-Based Learning} / PBL):}
    Model pembelajaran yang di dalamnya siswa berkolaborasi untuk menghasilkan artefak proyek nyata \autocite{Zhang2023}.
    \textbf{Analisis terhadap \textit{Over-personalization}:} Sangat baik untuk kolaborasi mendalam, namun memiliki kendala besar dalam penilaian otomatis. Sulit bagi sistem untuk memisahkan kontribusi individu dari hasil kelompok (\textit{free-rider problem}) secara akurat tanpa intervensi manual dosen \autocite{Ryoo2021}, sehingga sulit dijadikan input data \textit{real-time} untuk algoritma adaptif.

    \item \textbf{Pendampingan Sejawat (\textit{Peer Mentoring}):}
    Mekanisme bimbingan \textit{one-on-one} antara siswa senior dan junior \autocite{Le2024}.
    \textbf{Analisis terhadap \textit{Over-personalization}:} Hubungan ini sering kali bersifat hierarkis dan searah (senior mengajar junior), bukan kolaborasi setara. Hal ini kurang efektif untuk tujuan memecahkan \textit{echo chamber} yang membutuhkan paparan terhadap berbagai perspektif sejawat yang beragam, bukan sekadar transfer pengetahuan vertikal.

    \item \textbf{Formasi Kelompok Berbasis Klastering (\textit{Group Formation}):}
    Menggunakan algoritma seperti K-Means untuk membentuk kelompok heterogen di awal pembelajaran \autocite{Nalli2022}.
    \textbf{Analisis terhadap \textit{Over-personalization}:} Efektif untuk "memaksa" keberagaman di awal (\textit{cold start}). Namun, karena sifatnya statis, teknik ini tidak menghasilkan data kinerja kontinu yang bisa digunakan sistem untuk memantau apakah siswa kembali mengisolasi diri setelah kelompok terbentuk.

    \item \textbf{Sesi Diskusi Terbuka (\textit{Open Discussion}):}
    Memfasilitasi forum bebas untuk pertukaran ide \autocite{DaifAllah2016}.
    \textbf{Analisis terhadap \textit{Over-personalization}:} Dalam forum bebas, siswa cenderung mencari teman yang sepemikiran (\textit{homophily}), yang justru memperkuat \textit{echo chamber} alih-alih memecahkannya. Selain itu, sulit menilai kualitas interaksi secara otomatis tanpa NLP yang kompleks.

    \item \textbf{Penilaian Sejawat (\textit{Peer Assessment}):}
    Siswa menilai pekerjaan rekan lain berdasarkan rubrik dengan status yang setara \autocite{Ihichr2024}.
    \textbf{Analisis terhadap \textit{Over-personalization}:} Teknik ini paling potensial karena sistem memegang kendali penuh atas "siapa menilai siapa". Sistem dapat menugaskan siswa untuk menilai rekan yang memiliki sudut pandang bertentangan atau tingkat kemampuan berbeda (di luar \textit{bubble}-nya). Selain itu, kualitas umpan balik (R3) dapat diukur sebagai metrik keterbukaan siswa terhadap kolaborasi \autocite{Lu2012}.
\end{itemize}

\subsubsection{Alternatif Teknik Pencegahan \textit{Echo Chamber} dan \textit{Filter Bubble}}
\label{sssec:alternatif_pencegahan_bubble}
Subbab ini mengevaluasi pendekatan spesifik yang telah dikembangkan dalam literatur sistem rekomendasi modern (Bab \ref{chap:studi-literatur}) untuk mengatasi masalah isolasi, guna melihat kesesuaiannya untuk diadopsi ke dalam sistem asesmen pendidikan.

\begin{itemize}
    \item \textbf{\textit{User-Controllable Recommender Systems} (UCRS):}
    Pendekatan yang memberikan fitur kendali eksplisit kepada pengguna untuk mengatur tingkat diversifikasi rekomendasi mereka sendiri \autocite{Wang2022}.
    \textbf{Analisis:} Pendekatan ini sangat bergantung pada kesadaran pengguna. Dalam konteks pendidikan, siswa sering kali tidak menyadari bahwa mereka berada dalam \textit{bubble} kompetensi ("\textit{unconscious incompetence}"). Mengandalkan siswa untuk secara manual meminta materi yang "tidak mereka sukai" berisiko tidak efektif dibandingkan intervensi otomatis.

    \item \textbf{\textit{Group Recommender Systems} (GRS):}
    Teknik yang merekomendasikan konten berdasarkan preferensi agregat dari sekelompok pengguna untuk meningkatkan keadilan (\textit{fairness}) dan pluralitas \autocite{Dokoupil2022}.
    \textbf{Analisis:} GRS cenderung memprioritaskan konsensus kelompok atau rata-rata preferensi. Dalam pembelajaran adaptif, hal ini berisiko mengorbankan personalisasi spesifik yang dibutuhkan siswa \textit{outlier} (sangat mahir atau sangat tertinggal), karena rekomendasi akan tertarik ke arah "pusat" kelompok, yang bisa menyebabkan \textit{under-personalization}.

    \item \textbf{Pemodelan Komunitas Implisit (\textit{FRediECH}):}
    Sebuah teknik rekomendasi teman yang dirancang khusus untuk memecahkan \textit{echo chamber} dengan merekomendasikan koneksi sosial yang beragam namun relevan, berdasarkan pembelajaran representasi interaksi pengguna \autocite{Tommasel2021}.
    \textbf{Analisis:} Teknik ini sangat kuat untuk memecah polarisasi opini dengan memperluas jejaring sosial siswa. Namun, penerapannya dalam asesmen kompetensi memiliki tantangan: FRediECH berfokus pada rekomendasi "siapa yang harus dijadikan teman", bukan "materi apa yang harus dipelajari". Mengadopsinya memerlukan modifikasi agar rekomendasi koneksi tersebut dapat diterjemahkan menjadi aktivitas belajar yang terukur (misalnya, menjadi mitra \textit{peer assessment}), bukan sekadar jejaring sosial pasif.
\end{itemize}

\subsection{Analisis Penentuan Solusi}
\label{ssec:analisis_penentuan_solusi}
Penentuan solusi teknis dilakukan menggunakan metode \textit{Weighted Decision Matrix} (Matriks Keputusan Terbobot). Pendekatan ini digunakan untuk memilih komponen sistem secara objektif berdasarkan kriteria yang diturunkan dari Batasan Masalah (Bab \ref{chap:pendahuluan}) dan Tinjauan Literatur (Bab \ref{chap:studi-literatur}).

Analisis dibagi menjadi tiga keputusan terpisah untuk memilih: (1) Teknik Asesmen Individu, (2) Teknik Interaksi Sosial, dan (3) Teknik Mitigasi \textit{Over-personalization}.

\subsubsection{Keputusan 1: Penentuan Teknik Asesmen Individu}
\label{sssec:keputusan_asesmen_individu}
Tujuannya adalah memilih algoritma pelacakan pengetahuan. Berdasarkan tinjauan literatur, kandidat utamanya adalah IRT, FA, BKT, DKT, dan Elo-rating.

\begin{longtable}{@{\extracolsep{\fill}} p{3.0cm} c c c c c c}
\caption{Matriks Keputusan: Teknik Asesmen Individu}\label{tbl:solusi-r1} \\
\toprule
\textbf{Kriteria} & \textbf{Bobot} & \textbf{IRT} & \textbf{FA} & \textbf{BKT} & \textbf{DKT} & \textbf{Elo} \\
\midrule
\endfirsthead
\caption*{Tabel \ref{tbl:solusi-r1} (Lanjutan)} \\
\toprule
\textbf{Kriteria} & \textbf{Bobot} & \textbf{IRT} & \textbf{FA} & \textbf{BKT} & \textbf{DKT} & \textbf{Elo} \\
\midrule
\endhead
\midrule
\multicolumn{6}{r}{\textit{Bersambung ke halaman berikutnya}} \\
\bottomrule
\endfoot
\bottomrule
\endlastfoot

% Data Penilaian
Efisiensi Komputasi & 30\% & 2 & 1 & 3 & 1 & 5 \\
Kebutuhan Data Awal (Cold Start) & 30\% & 2 & 2 & 3 & 2 & 5 \\
Akurasi Pelacakan & 20\% & 4 & 3 & 3 & 5 & 4 \\
Interpretabilitas & 20\% & 4 & 3 & 4 & 1 & 5 \\
\midrule
\textbf{Skor Terbobot} & \textbf{100\%} & \textbf{2.8} & \textbf{2.1} & \textbf{3.2} & \textbf{2.1} & \textbf{4.8} \\
\end{longtable}

\textbf{Rasionalisasi Objektif dan Analisis \textit{Trade-off}:}
\begin{itemize}
    \item \textbf{Efisiensi \& Kebutuhan Data (Bobot Tinggi - 60\%):} Kriteria ini diberi bobot tertinggi mengingat batasan masalah pengembangan purwarupa yang mungkin tidak memiliki dataset historis besar. \textbf{Elo-rating} unggul mutlak (skor 5) karena algoritma ini dirancang untuk bekerja tanpa pra-kalibrasi yang rumit dan efisien secara komputasi \autocite{Vesin2022}. Sebaliknya, \textbf{IRT} dan \textbf{FA} membutuhkan sampel besar untuk kalibrasi parameter butir soal, dan \textbf{DKT} membutuhkan data pelatihan masif untuk melatih jaringan saraf \autocite{Minn2022}.
    
    \item \textbf{Akurasi vs. Kelayakan (Trade-off):} \textbf{DKT} memiliki skor akurasi tertinggi (5) karena kemampuannya menangkap pola kompleks \autocite{Minn2022}. Namun, terjadi \textit{trade-off} signifikan pada aspek \textbf{Interpretabilitas} (skor 1) karena sifat \textit{black-box}-nya, yang menyulitkan sistem untuk menjelaskan alasan rekomendasi kepada siswa. \textbf{Elo-rating} dipilih bukan karena ia paling akurat, tetapi karena ia adalah solusi paling layak (\textit{feasible}) untuk diimplementasikan dalam skenario data terbatas dengan akurasi yang masih dapat diterima.
\end{itemize}

\subsubsection{Keputusan 2: Penentuan Teknik Interaksi Sosial}
\label{sssec:keputusan_asesmen_sosial}
Tujuannya adalah memilih bentuk aktivitas sosial yang dapat memfasilitasi kolaborasi sekaligus menghasilkan data yang dapat diolah sistem.

\begin{longtable}{@{\extracolsep{\fill}} p{3.0cm} c c c c c c}
\caption{Matriks Keputusan: Teknik Interaksi Sosial}\label{tbl:solusi-r2} \\
\toprule
\textbf{Kriteria} & \textbf{Bobot} & \textbf{PBL} & \textbf{Mentor} & \textbf{Group} & \textbf{Diskusi} & \textbf{Peer Assmt.} \\
\midrule
\endfirsthead
\caption*{Tabel \ref{tbl:solusi-r2} (Lanjutan)} \\
\toprule
\textbf{Kriteria} & \textbf{Bobot} & \textbf{PBL} & \textbf{Mentor} & \textbf{Group} & \textbf{Diskusi} & \textbf{Peer Assmt.} \\
\midrule
\endhead
\midrule
\multicolumn{6}{r}{\textit{Bersambung ke halaman berikutnya}} \\
\bottomrule
\endfoot
\bottomrule
\endlastfoot

% Data Penilaian
Keterukuran Data (Measurability) & 35\% & 2 & 3 & 1 & 1 & 5 \\
Kesetaraan Interaksi & 25\% & 5 & 2 & 5 & 5 & 5 \\
Potensi Eksposur Keberagaman & 25\% & 4 & 2 & 3 & 3 & 5 \\
Kompleksitas Implementasi & 15\% & 2 & 3 & 4 & 3 & 4 \\
\midrule
\textbf{Skor Terbobot} & \textbf{100\%} & \textbf{3.25} & \textbf{2.5} & \textbf{3.0} & \textbf{2.8} & \textbf{4.85} \\
\end{longtable}

\textbf{Rasionalisasi Objektif dan Analisis \textit{Trade-off}:}
\begin{itemize}
    \item \textbf{Keterukuran Data (35\%):} Sistem membutuhkan input kuantitatif untuk algoritma (FR-06). \textbf{Peer Assessment} mendapatkan skor tertinggi (5) karena menghasilkan artefak terstruktur (skor rubrik dan teks umpan balik) yang mudah dikuantifikasi \autocite{Ihichr2024}. Sebaliknya, \textbf{Diskusi Terbuka} dan \textbf{PBL} (Project-Based Learning) sangat sulit dinilai kontribusi individunya secara otomatis tanpa intervensi manual dosen \autocite{Ryoo2021}.
    
    \item \textbf{Kesetaraan \& Eksposur (Trade-off):} \textbf{Mentoring} mendapat skor rendah pada kesetaraan (2) karena sifatnya hierarkis (senior-junior) \autocite{Le2024}, yang kurang cocok untuk memecah \textit{echo chamber} dibandingkan interaksi setara. \textbf{PBL} sebenarnya sangat baik dalam kesetaraan dan kolaborasi mendalam (skor 5), namun dipilihnya \textbf{Peer Assessment} merupakan \textit{trade-off} logis karena PBL memiliki kompleksitas implementasi yang tinggi dan risiko \textit{free-rider} yang sulit dideteksi sistem \autocite{Zhang2023}. Peer Assessment dipilih karena memungkinkan sistem "mengatur" siapa menilai siapa untuk memaksa eksposur keberagaman.
\end{itemize}

\subsubsection{Keputusan 3: Penentuan Teknik Mitigasi \textit{Over-personalization}}
\label{sssec:keputusan_mitigasi_bubble}
Tujuannya adalah memilih pendekatan algoritmik untuk mengatasi isolasi pembelajaran, berdasarkan alternatif yang tersedia di Bab II.

\begin{longtable}{@{\extracolsep{\fill}} p{3.5cm} c c c c}
\caption{Matriks Keputusan: Teknik Mitigasi Bubble}\label{tbl:solusi-r3} \\
\toprule
\textbf{Kriteria} & \textbf{Bobot} & \textbf{Alt A: UCRS} & \textbf{Alt B: GRS} & \textbf{Alt C: FRediECH} \\
\midrule
\endfirsthead
\caption*{Tabel \ref{tbl:solusi-r3} (Lanjutan)} \\
\toprule
\textbf{Kriteria} & \textbf{Bobot} & \textbf{Alt A: UCRS} & \textbf{Alt B: GRS} & \textbf{Alt C: FRediECH} \\
\midrule
\endhead
\midrule
\multicolumn{5}{r}{\textit{Bersambung ke halaman berikutnya}} \\
\bottomrule
\endfoot
\bottomrule
\endlastfoot

% Data Penilaian
Otomasi Intervensi & 40\% & 1 & 4 & 4 \\
Preservasi Personalisasi Individu & 30\% & 5 & 2 & 3 \\
Relevansi Pemecahan Echo Chamber & 30\% & 3 & 3 & 5 \\
\midrule
\textbf{Skor Terbobot} & \textbf{100\%} & \textbf{2.8} & \textbf{3.1} & \textbf{4.0} \\
\end{longtable}

\textbf{Rasionalisasi Objektif dan Analisis \textit{Trade-off}:}
\begin{itemize}
    \item \textbf{Otomasi Intervensi (40\%):} Untuk mengatasi siswa yang tidak sadar akan inkompetensinya (\textit{unconscious incompetence}), sistem harus proaktif. \textbf{UCRS} (User-Controllable) mendapat skor terendah (1) karena sangat bergantung pada inisiatif siswa untuk mengubah pengaturan rekomendasi \autocite{Wang2022}. \textbf{GRS} dan \textbf{FRediECH} unggul karena bekerja secara otomatis di latar belakang.
    
    \item \textbf{Preservasi Personalisasi (30\%):} \textbf{GRS} (Group Recommender) cenderung merekomendasikan apa yang disukai mayoritas kelompok, yang berisiko mengorbankan kebutuhan unik individu (skor 2) \autocite{Dokoupil2022}. \textbf{UCRS} paling menghormati otonomi individu (skor 5). \textbf{FRediECH} berada di tengah-tengah; ia merekomendasikan koneksi orang lain, namun tetap mempertahankan profil individu pengguna.
    
    \item \textbf{Relevansi Pemecahan Echo Chamber (30\%):} \textbf{FRediECH} secara spesifik dirancang dengan fungsi tujuan (\textit{objective function}) untuk memecahkan \textit{echo chamber} dengan merekomendasikan teman yang beragam \autocite{Tommasel2021}. Ini lebih relevan dibanding GRS yang fokus pada keadilan kelompok.
    
    \item \textbf{Kesimpulan:} \textbf{FRediECH (Implisit Community Modeling)} dipilih sebagai pendekatan mitigasi. Meskipun UCRS memberikan kontrol lebih besar, risiko pasifnya siswa membuat FRediECH lebih unggul untuk tujuan penelitian ini. Konsep FRediECH akan diadaptasi: alih-alih sekadar merekomendasikan "teman", sistem akan merekomendasikan "mitra interaksi" (untuk aktivitas Peer Assessment yang terpilih di Keputusan 2) guna memecah isolasi.
\end{itemize}

\subsubsection{Ringkasan Solusi Terpilih}
Berdasarkan analisis matriks keputusan di atas, purwarupa sistem akan mengimplementasikan kombinasi berikut:
\begin{enumerate}
    \item \textbf{Asesmen Individu:} Menggunakan \textbf{\textit{Elo-rating}} karena efisiensinya dalam kondisi data terbatas.
    \item \textbf{Interaksi Sosial:} Menggunakan \textbf{\textit{Peer Assessment}} karena menghasilkan data kuantitatif yang terukur.
    \item \textbf{Mitigasi Isolasi:} Mengadaptasi logika \textbf{FRediECH} untuk merekomendasikan mitra penilaian sejawat yang berada di luar \textit{bubble} pengguna, guna memastikan terjadinya eksposur terhadap keberagaman.
\end{enumerate}